\subsection{Qué hacía el TP x86}

El TP3 original arrancaba en modo real desde el boot-sector, habilitaba A20, cargaba la GDT con \texttt{lgdt}, prendía PE de CR0,
hacía el \texttt{jmp} lejano a modo protegido, cargaba los segmentos y fijaba el stack en \texttt{0x27000}. Además declaraba
un segmento de pantalla aparte y limpiaba la VGA.

\subsection{Equivalente en RISC-V}

En RISC-V no hay segmentación ni A20. Los permisos de memoria salen de la PTE Sv39 (\texttt{V}, \texttt{R}, \texttt{W},
\texttt{X}, \texttt{U}, \texttt{G}, \texttt{A}, \texttt{D}), y lo más parecido a ``pasar a modo protegido'' es
\textbf{ejecutar \texttt{mret} desde M-mode con \texttt{mstatus.MPP}$=$1 (S-mode) y
\texttt{mepc}$=$\texttt{kernel\_main}}.

La cadena de arranque es:

\begin{enumerate}
  \item QEMU salta a \texttt{\_start} en \texttt{0x80000000} (M-mode).
  \item \texttt{entry.S} inicializa \texttt{gp}, stack del kernel y limpia \texttt{.bss}, y llama a
  \texttt{machine\_init}.
  \item \texttt{machine.S::machine\_init} (una ``mini firmware'' en M-mode, ver sección~\ref{sec:machine}) configura
  PMP, delega traps a S-mode, programa el tick y ejecuta \texttt{mret} $\to$ \texttt{kernel\_main} ya en S-mode.
\end{enumerate}

\subsection{\texttt{riscv/entry.S}}

Fragmento relevante del entrypoint:

\begin{codesnippet}
\begin{verbatim}
 _start:
     .option push
     .option norelax
     la gp, __global_pointer$
     .option pop

     csrw mie, zero
     csrci mstatus, 0x8        /* MIE */
     csrw sie, zero
     csrci sstatus, 0x2        /* SIE */

     la sp, kernel_stack_top

     /* Clear .bss */
     la t0, __bss_start
     la t1, __bss_end
 1:
     bgeu t0, t1, 2f
     sb zero, 0(t0)
     addi t0, t0, 1
     j 1b
 2:
     call machine_init
\end{verbatim}
\end{codesnippet}

Detalles técnicos:

\begin{itemize}
  \item \textbf{\texttt{gp} y \texttt{.option norelax}.} La psABI reserva \texttt{gp} para accesos \emph{small-data}
  gp-relativos~\cite{riscv-psabi}. Si no se desactiva la \emph{linker relaxation} alrededor de \texttt{la gp,
  \_\_global\_pointer\$}, el linker ve un \texttt{auipc}/\texttt{addi} y lo reescribe como un \texttt{addi gp, gp, 0}
  gp-relativo --- básicamente, asume que \texttt{gp} ya vale lo que tiene que valer. Resultado: \texttt{gp} queda en
  basura y cualquier acceso small-data posterior lee cualquier cosa. \texttt{.option norelax} apaga esa optimización en
  ese bloque puntual.

  \item \textbf{Stack del kernel.} Es el equivalente a ``setear el stack en \texttt{0x27000}'' del TP x86. El símbolo
  \texttt{kernel\_stack\_top} está definido al final de un bloque en \texttt{.bss} de 8~KiB, resultado del linker (ver
  \texttt{riscv/linker.ld}). No hace falta una dirección mágica porque el linker coloca todo a partir de
  \texttt{0x80000000} y la ubicación concreta de el stack la decide el linker al resolver símbolos.
\end{itemize}

\subsection{\texttt{riscv/machine.S}: mini firmware en M-mode}
\label{sec:machine}

\texttt{machine\_init} cumple el rol que en un sistema típico cumple OpenSBI (la firmware SBI de referencia en
RISC-V~\cite{riscv-sbi}). Como se corre con \texttt{-bios none}, no hay SBI, y el kernel tiene que proveer él mismo la
configuración mínima de M-mode. El esqueleto sigue la receta \emph{Bare Bones} para RISC-V documentada en
OSDev~\cite{osdev-riscv}. Los pasos:

\begin{enumerate}
  \item \textbf{Stack y trap vector de M-mode.} \texttt{mscratch} guarda el tope de el stack de M-mode y \texttt{mtvec}
  apunta a \texttt{machine\_trap\_entry}.
  \item \textbf{Delegación.} \texttt{mideleg} delega a S-mode los interrupts de supervisor (SSIP, STIP, SEIP) y
  \texttt{medeleg} delega las excepciones estándar 0..15. Así, las excepciones desde U-mode saltan directamente a
  \texttt{stvec} del kernel y no a \texttt{mtvec}.
  \item \textbf{PMP.} Una sola región TOR que abarca todo el espacio (\texttt{pmpaddr0 = -1}, \texttt{pmpcfg0 =
  RWX+TOR}) le da a S-mode y U-mode acceso pleno según PMP. El control fino después lo dan las PTEs Sv39.
  \item \textbf{Timer.} Se habilita MTIE en \texttt{mie} y se programa el primer \texttt{mtimecmp} = \texttt{time} +
  \texttt{g\_timer\_delta\_ticks}.
  \item \textbf{Salto a S-mode.} Se setea \texttt{mstatus.MPP}$=$1 (S), \texttt{mstatus.MPIE}$=$1,
  \texttt{mepc}$=$\texttt{kernel\_main} y se ejecuta \texttt{mret}. Esto es el análogo RISC-V del \texttt{jmp
  0xA0:modo\_protegido} del TP x86.
\end{enumerate}

\begin{codesnippet}
\begin{verbatim}
 # Delegar interrupts a S-mode: SSIP, STIP, SEIP
 li t0, ((1 << 1) | (1 << 5) | (1 << 9))
 csrw mideleg, t0

 # Delegar excepciones 0..15 a S-mode
 li t0, 0xFFFF
 csrw medeleg, t0

 # PMP: R|W|X + TOR cubriendo 0..max
 li t0, -1
 csrw pmpaddr0, t0
 li t0, 0x0F
 csrw pmpcfg0, t0
\end{verbatim}
\end{codesnippet}

\subsection{``Segmento de pantalla'' y limpieza inicial}

Como no hay segmentación, la pantalla se resuelve por \textbf{mapeo de MMIO}: el UART 16550 de QEMU \texttt{virt} en
\texttt{0x1000\_0000}. En cada page-table de tarea (y en la de idle) se mapea ese rango con una megapage (2~MiB) a nivel
L1, con flags \texttt{V|R|W|A|D}, sin \texttt{U}, para que el kernel pueda escribirlo. El equivalente al ``segmento de
pantalla sólo accesible por el kernel'' del TP x86 es entonces el bit \texttt{U}$=$0 de esa entrada.

La limpieza y pintado del tablero equivalen, a nivel código, a \texttt{screen\_draw} / \texttt{screen\_drawBox} en
\texttt{riscv/screen.cpp}. Internamente la pantalla se emula como un buffer \texttt{g\_video[50][80]} de celdas
(caracter + atributo estilo VGA), que se vuelca al UART mediante secuencias ANSI (\texttt{\textbackslash x1b[?1049h}
para el alternate screen, \texttt{\textbackslash x1b[y;xH} para posicionar el cursor, \texttt{\textbackslash x1b[fg;bgm}
para atributos, etc.). Al imprimir, un atributo estilo \texttt{C\_BG\_RED | C\_FG\_WHITE} se traduce a los códigos ANSI
\texttt{41} y \texttt{97}.
